\section{Interpretation and limits}
There are two distinct ways that more recent observations enter a forecast. An additive estimator replaces older increments with newer unconditional event fractions. A multiplicative estimator also changes the base on which the next growth factor operates. When a cohort has had very few early charge-offs, a sequence of estimated growth factors can propagate a small change in early incidence into a large change in the projected count. Comparing two information windows within each family exposes this difference. The ordering of the two windows in a given historical replay therefore has limited value without the development formula and the individual vintage errors.

Two checks added after the primary comparison address selection and calibration. Retaining all nonnegative approval-to-disbursement delays increases the scored sample to 918,448 \sba{} and 105,171 504 loans. Period versus mature additive MAE is then 5.02 versus 6.26 pp for \sba{} and 3.17 versus 3.99 pp for 504. Moving from the mature development factor to the age-to-age chain still increases MAE in both programs. The period upper forecast again covers zero of seventeen overlapping \sba{} vintages and seven of seventeen 504 vintages. The original entry restriction therefore has a substantial effect on 504 portfolio size while preserving these directional findings.

Keeping the original cohorts and calibrating a separate historical margin for each point forecast changes the picture. During the crisis overlap the mature development factor's upper forecast covers thirteen of seventeen \sba{} outcomes, compared with zero for the period additive center and eleven for the age-to-age chain. In 504 the corresponding counts are eight, seven and eleven. Thus the zero-of-seventeen result is specific to the combination of a period increment forecast and its lagged error margin. Changing the forecast center changes the errors available for calibration as well as the point estimate.

A second operational issue is the lag in feedback. At an eight-quarter landmark the errors used to set the margin come from vintages at least twelve quarters older. The newest errors in a calibration sample can still describe conditions that have passed. This lag is fixed by the outcome horizon even when recent age-specific observations can already influence the point forecast. An archive-based report can make the lag visible by publishing the issue date, most recent mature cohort and largest subsequent breach next to the stated target coverage.

The historical exercise covers one US financial crisis. Adjacent portfolios and the two SBA programs share national economic conditions. The small number of archived issue dates further limits inference about publication revisions in other periods. The upper forecasts measure empirical shortfall risk for a counted portfolio. They have no exchangeability guarantee under a common shock. The public fields identify administrative recognition dates. Changes in collections, accounting and reporting can move them even if underlying borrower distress has a different time path. The entry rule selects relatively prompt disbursements and the analysis gives each loan the same weight. The measured errors apply to these public portfolios. An underwriting, reserve-dollar or macroeconomic capital application would need its own loss definition, exposure weights and validation data.
